% Appendix D: Generative AI Usage Declaration — Handbook 9.7
\section{Generative AI Usage}
\label{app:ai}

This appendix declares the use of generative AI tools during the project,
in accordance with the department's policy on generative AI (Handbook
Section~9.7).

\subsection*{Tools Used}

\begin{table}[H]
\centering
\small
\renewcommand{\arraystretch}{1.3}
\begin{tabular}{|l|l|}
\hline
\textbf{Tool} & \textbf{Provider} \\
\hline
Claude Code (Opus 4) & Anthropic \\
\hline
\end{tabular}
\caption{Generative AI tools used.}
\end{table}

\subsection*{How AI Was Used}

The following describes how Claude Code was used during the project, in
chronological order of development.  In each case the prompt was given
within the context of the full project codebase.  All core application logic (auction coordination, MPC orchestration,
Veilid integration, cryptographic protocols) was designed and written by
the author.  Claude Code was used primarily for infrastructure and
tooling (devnet setup, build scripts, UI), refactoring, and debugging.
All code, whether AI-generated or hand-written, was reviewed, tested,
and committed by the author.  All prose in this report was written by
the author; Claude Code was used only for cross-referencing and gap
analysis.

\begin{enumerate}
  % --- Nov 2025: project setup ---
  \item \textbf{Docker devnet bootstrap.}\\
    \textit{Prompt:} ``The Veilid devnet merge request doesn't work.
    Nodes fail to bootstrap.  Here's the docker-compose and the error
    output: [pasted logs].''\\
    \textit{Output:} Claude iterated on the Docker Compose configuration,
    bootstrap entrypoint scripts, network key settings, and DNS
    configuration.\\
    \textit{What was done:} The author tested each iteration against the
    actual Veilid bootstrap process, feeding error logs back across
    multiple rounds until the 20-node devnet started reliably.

  % --- Jan 2026: core infrastructure ---
  \item \textbf{IP spoofing LD\_PRELOAD library.}\\
    \textit{Prompt:} ``Write a C LD\_PRELOAD library that intercepts
    \texttt{bind()}, \texttt{connect()}, \texttt{sendto()}, and
    \texttt{recvfrom()} to translate fake public IPs (1.2.3.x) to
    localhost, so multiple Veilid nodes can run on one machine.''\\
    \textit{Output:} Claude generated \texttt{ip\_spoof.c} with the
    socket interception logic and address translation table.\\
    \textit{What was done:} The author reviewed the implementation,
    tested it against the devnet, and later extended it with NAT table
    support for Tailscale.  The library was subsequently ported to Rust
    by Claude for the veilid-playground project.

  \item \textbf{Connecting market nodes to the devnet.}\\
    \textit{Prompt:} ``Now adapt the market application's Veilid node
    config to connect to this dockerised devnet.''\\
    \textit{Output:} Claude generated the node configuration changes:
    network key, bootstrap address, port offsets, and capability flags.\\
    \textit{What was done:} The author tested the connection against the
    running devnet and iterated on capability settings.

  \item \textbf{Dioxus UI: initial layout.}\\
    \textit{Prompt:} ``Build a Dioxus desktop UI with three panels:
    marketplace listings, active auctions, and a bid submission form.
    Use the AppState struct in src/app/state.rs.''\\
    \textit{Output:} Claude generated the Dioxus component hierarchy
    (\texttt{src/app/components.rs}, \texttt{src/app/mod.rs}) with
    state bindings.\\
    \textit{What was done:} The author reviewed the generated code,
    adjusted the layout, and integrated it with the Veilid node
    lifecycle.

  \item \textbf{Dioxus UI: stylesheet refactor.}\\
    \textit{Prompt:} ``Refactor the inline styles into a Dioxus
    stylesheet.''\\
    \textit{Output:} Claude extracted all inline styles into a
    centralised stylesheet.\\
    \textit{What was done:} The author reviewed and committed the
    refactor.

  \item \textbf{Build scripting: 3-node demo.}\\
    \textit{Prompt:} ``Write a Makefile target that starts a 3-node demo:
    builds the release binary, starts the devnet, then launches 3 market
    nodes with offsets 20--22 using LD\_PRELOAD.''\\
    \textit{Output:} Claude generated the \texttt{make demo} target
    with process management and SIGTERM cleanup.\\
    \textit{What was done:} The author tested the target, fixed a
    platform-specific issue on Linux, and committed it.

  % --- Feb 2026: architecture cleanup ---
  \item \textbf{Unit and integration test harness.}\\
    \textit{Prompt:} ``Set up a mock-based test infrastructure so I can
    simulate $N$-party auctions without a real devnet.''\\
    \textit{Output:} Claude generated \texttt{MockDht},
    \texttt{MockTime}, \texttt{MockRandom}, and the
    \texttt{MultiPartyHarness} for running full auction flows in
    under 1\,s.\\
    \textit{What was done:} The author reviewed the mock implementations,
    verified they matched real Veilid DHT semantics, and built
    integration tests on top of them.

  \item \textbf{Two-tier registry refactor.}\\
    \textit{Prompt:} ``The shared DHT registry lets any party delete
    others' listings.  Refactor to a two-tier registry-of-registries
    pattern with per-seller owned records.''\\
    \textit{Output:} Claude restructured the registry into per-seller
    DHT records with a top-level catalog.\\
    \textit{What was done:} The author verified the permission model
    and tested multi-seller scenarios.

  \item \textbf{Codebase cleanup and refactoring.}\\
    \textit{Prompt:} ``Remove all dead code, collapse unnecessary
    generic parameters, deduplicate test helpers, and get clippy passing
    with pedantic lints.''\\
    \textit{Output:} Claude produced a series of refactoring patches:
    collapsing \texttt{AuctionLogic<D,T,M,C>} to \texttt{<D,C>},
    deleting unused traits and mock self-tests, renaming
    \texttt{MpcSidecar} to \texttt{MpcTunnelProxy}, deduplicating
    broadcast helpers, unifying \texttt{MockDht} types, and dropping
    \texttt{async-std} in favour of Tokio.  Net result: $-$2,877 lines
    across six commits.\\
    \textit{What was done:} The author reviewed each patch, verified
    all 169 unit and 63 integration tests still passed, and committed
    as six logical commits.

  \item \textbf{Error hardening.}\\
    \textit{Prompt:} ``Audit the codebase for panics and unwraps in
    production code paths.  Replace with proper error handling.''\\
    \textit{Output:} Claude identified unsafe \texttt{unwrap()} calls
    in node configuration, string indexing, and insecure storage setup,
    and generated replacements using \texttt{Result} propagation.\\
    \textit{What was done:} The author reviewed each change, confirmed
    no behaviour regressions, and committed the fixes.

  \item \textbf{Strict clippy lints.}\\
    \textit{Prompt:} ``Enable pedantic, nursery, and cargo clippy lint
    groups and fix everything.''\\
    \textit{Output:} Claude enabled the lint groups and fixed the
    resulting warnings across the codebase.\\
    \textit{What was done:} The author reviewed the changes and
    committed them.

  % --- Feb 2026: E2E and MPC ---
  \item \textbf{E2E test debugging: sequential auctions.}\\
    \textit{Prompt:} ``All 3 E2E tests are failing with route timeout on
    the second auction.  The first auction completes fine.  Here's the log
    output: [pasted test log].''\\
    \textit{Output:} Claude identified that broadcast routes from the
    first auction were stale by the time the second auction started, and
    suggested adding a route refresh after MPC completion.\\
    \textit{What was done:} The author implemented the
    \texttt{needs\_route\_refresh} AtomicBool flag in
    \texttt{MpcOrchestrator} based on this diagnosis, then verified the
    fix resolved the failure.

  \item \textbf{E2E test debugging: concurrent auction deadlock.}\\
    \textit{Prompt:} ``Two auctions expiring at the same time cause a
    deadlock.  Here's the log: [pasted log].''\\
    \textit{Output:} Claude identified that nodes processed expired
    listings in different orders, and suggested sorting them
    deterministically by key before sequential processing.\\
    \textit{What was done:} The author implemented the deterministic
    sort and verified the deadlock was resolved.

  \item \textbf{Performance optimisation: MPC tunnel pipelining.}\\
    \textit{Prompt:} ``E2E tests take 900s.  Profile the MPC tunnel
    and find the bottleneck.''\\
    \textit{Output:} Claude suggested pipelining \texttt{app\_call}
    sends with a bounded semaphore (depth 8) instead of sequential
    send-wait cycles, and reducing polling intervals.\\
    \textit{What was done:} The author implemented the pipeline in
    \texttt{MpcTunnelProxy}, tuned the semaphore depth experimentally,
    and measured a reduction from ${\sim}$900\,s to ${\sim}$570\,s.

  % --- Mar 2026: UI polish ---
  \item \textbf{Dioxus UI: auction progress indicators.}\\
    \textit{Prompt:} ``Add visual feedback for the MPC phases: route
    collection, readiness barrier, and MPC execution.''\\
    \textit{Output:} Claude generated status indicators showing each
    phase's progress in the auction view.\\
    \textit{What was done:} The author wired the indicators to the
    orchestrator state callbacks and tested them.

  \item \textbf{Dioxus UI: live hex traffic panel.}\\
    \textit{Prompt:} ``Add a live hexdump display showing the encrypted
    MPC traffic flowing through the tunnel.  Make it look nice.''\\
    \textit{Output:} Claude generated a formatted hex viewer component
    with colour-coded bytes and ASCII sidebar.\\
    \textit{What was done:} The author integrated it into the auction
    detail view.

  \item \textbf{Dioxus UI: dark theme and 3D lattice header.}\\
    \textit{Prompt:} ``Restyle the UI with a dark theme and add a
    rotating 3D lattice animation for the header.''\\
    \textit{Output:} Claude generated the dark-themed CSS and a WebGL
    lattice animation.\\
    \textit{What was done:} The author reviewed and committed the
    styling changes.

  \item \textbf{Dioxus UI: graceful shutdown.}\\
    \textit{Prompt:} ``Handle SIGTERM so the UI and Veilid node shut
    down cleanly instead of hanging.''\\
    \textit{Output:} Claude implemented signal handling with Tokio's
    signal listener and a clean shutdown sequence.\\
    \textit{What was done:} The author tested the shutdown behaviour
    with \texttt{make demo} and committed it.

  % --- Mar 2026: benchmarking ---
  \item \textbf{Benchmark harness.}\\
    \textit{Prompt:} ``Write a benchmark binary that runs automated
    multi-party auction trials across configurable party counts and
    devnet sizes, with CSV output and resume support.''\\
    \textit{Output:} Claude generated the \texttt{bench-auction} binary
    and associated shell scripts.\\
    \textit{What was done:} The author tested the harness across 40, 60,
    and 80-node devnets and committed it.

  \item \textbf{Veilid playground migration.}\\
    \textit{Prompt:} ``Migrate the E2E test infrastructure from Docker
    to the veilid-playground (native processes with LD\_PRELOAD).''\\
    \textit{Output:} Claude adapted the test helpers to detect and use
    the playground instead of Docker.\\
    \textit{What was done:} The author verified E2E tests passed on
    both Docker and playground backends.

  % --- Mar 2026: security audit ---
  \item \textbf{Security audit fixes.}\\
    \textit{Prompt:} ``Audit the codebase for security issues:
    race conditions, information leakage, serde misconfigurations.''\\
    \textit{Output:} Claude identified TOCTOU races in MPC session
    dedup and auction activation, serde annotation issues that could
    leak key/nonce material, and timing-leak vectors.\\
    \textit{What was done:} The author reviewed each finding, committed
    the fixes, and verified tests still passed.

  % --- Mar--Apr 2026: report ---
  \item \textbf{Benchmark data processing and graphing.}\\
    \textit{Prompt:} ``Write Python scripts to parse the benchmark CSV
    files, compute medians and success rates, and output .dat files for
    pgfplots.''\\
    \textit{Output:} Claude generated the data processing scripts that
    produced the \texttt{.dat} files used by the pgfplots charts in
    Chapter~5.\\
    \textit{What was done:} The author verified the computed values
    against the raw CSV data and adjusted chart formatting.

  \item \textbf{LaTeX formatting.}\\
    \textit{Prompt:} ``Set up the dissertation LaTeX build: XeLaTeX with
    Aptos font, Harvard-style biblatex, word count automation, pgfplots
    for benchmark charts.''\\
    \textit{Output:} Claude configured the document preamble, bibliography
    style, and pgfplots chart templates.\\
    \textit{What was done:} The author adjusted formatting to match
    handbook requirements and populated all content.

  \item \textbf{Report gap analysis.}\\
    \textit{Prompt:} ``Compare my Ch5 Results against the codebase and
    benchmark CSVs.  What's missing or inconsistent?''\\
    \textit{Output:} Claude flagged that several table values used
    approximate markers where exact data existed in the benchmark CSV,
    and that the warm/cold devnet distinction had no backing data.\\
    \textit{What was done:} The author replaced approximate values with
    exact figures from the CSV and removed the unsupported warm/cold
    comparison.
\end{enumerate}

\subsection*{Responsibility}

The author takes full responsibility for the content, accuracy, and
originality of this report and all project outputs.  AI tools were used as
assistants, not as authors.  All claims, analyses, and conclusions are the
author's own.
